\begin{figure*}[t]
\centering
\resizebox{\textwidth}{!}{%
\begin{tikzpicture}[
  font=\footnotesize, >=latex,
  strip/.style={draw=black!40, rounded corners=2pt, fill=black!3, align=center,
                inner sep=3pt, font=\scriptsize, minimum height=8mm},
  tower/.style={draw=black!55, rounded corners=3pt, fill=black!5,
                minimum width=16mm, minimum height=14mm, align=center, font=\scriptsize},
  op/.style={circle, draw=black!70, fill=black!8, inner sep=1.6pt, font=\normalsize},
  saebox/.style={draw=blue!45, rounded corners=3pt, fill=blue!5, align=center,
                 font=\scriptsize, minimum height=15mm, minimum width=26mm},
  acc/.style={draw=orange!80!black, fill=orange!20, rounded corners=2pt, align=center,
              font=\scriptsize},
  readout/.style={draw=black!55, rounded corners=2pt, fill=green!8, align=center,
                  inner sep=3pt, font=\scriptsize},
  a/.style={->, semithick, black!60},
  aa/.style={->, semithick, orange!85!black},
  st/.style={circle, fill=black!80, text=white, font=\scriptsize\bfseries, inner sep=1.3pt},
  lbl/.style={font=\scriptsize\itshape, black!55}
]

% ---- Stage 1: serialize ----
\node[strip] (logs) {insider-threat\\logs};
\node[strip, right=4mm of logs] (ser) {serialized user-day\\\texttt{DAY}/\texttt{PSY}/\texttt{SES}};
\draw[a] (logs) -- (ser);
\node[st, above=0.4mm of logs] {1};

% ---- Stage 2: base vs adapted towers + delta ----
\node[tower, above right=2mm and 11mm of ser] (base) {frozen\\base LM};
\node[tower, below right=2mm and 11mm of ser] (ada)  {adapted LM\\{\tiny(base\,+\,QLoRA)}};
\draw[a] (ser.east) -- (base.west);
\draw[a] (ser.east) -- (ada.west);
\node[op, right=17mm of ser] (sub) {$-$};
\draw[a] (base.east) -| (sub);
\draw[a] (ada.east)  -| (sub);
\node[acc, right=5mm of sub] (delta) {token delta\\$\delta_{\ell,t}$};
\draw[aa] (sub) -- (delta);
\node[st, above=1mm of base] {2};

% ---- Stage 3: top-k delta-SAE ----
\node[saebox, right=6mm of delta, minimum width=30mm, minimum height=17mm] (sae) {};
\draw[a] (delta) -- (sae);
\node[font=\scriptsize] at ([yshift=5mm]sae.center) {top-$k$ delta-SAE};
\node[font=\tiny, black!70] at ([xshift=-11mm]sae.center) {enc};
\node[font=\tiny, black!70] at ([xshift=11mm]sae.center)  {dec};
\draw[a] ([xshift=-8.5mm]sae.center) -- ([xshift=-2.5mm]sae.center);
\draw[a] ([xshift=2.5mm]sae.center)  -- ([xshift=8.5mm]sae.center);
% sparse-code glyph centered (few cells lit = top-k)
\begin{scope}[shift={(sae.center)}]
  \draw[draw=black!45, fill=black!10]  (-1mm, 3.2mm) rectangle ++(2mm,-1.3mm);
  \draw[draw=black!45, fill=black!10]  (-1mm, 1.7mm) rectangle ++(2mm,-1.3mm);
  \draw[draw=black!45, fill=orange!75] (-1mm, 0.2mm) rectangle ++(2mm,-1.3mm);
  \draw[draw=black!45, fill=black!10]  (-1mm,-1.3mm) rectangle ++(2mm,-1.3mm);
  \draw[draw=black!45, fill=orange!75] (-1mm,-2.8mm) rectangle ++(2mm,-1.3mm);
  \draw[draw=black!45, fill=black!10]  (-1mm,-4.3mm) rectangle ++(2mm,-1.3mm);
\end{scope}
\node[font=\tiny, black!55] at ([yshift=-6.7mm]sae.center) {code $z$ ($k$ active) $+\,\varepsilon$};
\node[st, above=1mm of sae] {3};

% ---- Stage 4: pick feature -> patch/ablate -> score effect ----
\node[acc, below=10mm of sae.south, xshift=-10mm] (feat) {gap-selected\\feature $f_j$};
\draw[aa] (sae.south) -- (feat.north);
\node[readout, right=6mm of feat, fill=orange!8, draw=orange!70!black] (inter)
  {patch $\tau$ / ablate $\nu$\\vs.\ matched control};
\draw[aa] (feat) -- (inter);
\node[readout, right=6mm of inter] (score) {$\Delta$ anomaly score\\$s \to s'$};
\draw[a] (inter) -- (score);
\node[st, above=0.4mm of feat] {4};

% ---- Stage 5: attribution ----
\node[readout, right=6mm of score] (attr) {attribution $m_f(c)$:\\profile vs.\ behavior};
\draw[a] (score) -- (attr);
\node[st, above=0.4mm of attr] {5};

\end{tikzpicture}}%
\caption{The delta-SAE audit, drawn as architecture.
\textbf{(1)}~A user-day is serialized into structured
\texttt{DAY}/\texttt{PSY}/\texttt{SES} text.
\textbf{(2)}~The same input runs through the frozen base model and its
benign-only QLoRA-adapted copy; at layer $\ell$ we take the per-token
residual-stream difference
$\delta_{\ell,t}=h^{\mathrm{ada}}_{\ell,t}-h^{\mathrm{base}}_{\ell,t}$
(the adapter is left abstract---the delta is the object of study).
\textbf{(3)}~$\delta$ feeds a top-$k$ sparse autoencoder (encoder $\to$ sparse
code $z$ with only $k$ units active, shown as the lit cells $\to$ decoder,
plus reconstruction error), whose decoder columns form the feature dictionary.
\textbf{(4)}~A gap-selected feature $f_j$ is causally patched ($\tau$) or
ablated ($\nu$) against a matched active control, and the effect on the adapted
anomaly score $s\!\to\!s'$ is measured.
\textbf{(5)}~Token attribution $m_f(c)$ labels the feature's content as
profile- or behavior-bound. A mechanistic claim requires both causal relevance
(Step~4) \emph{and} content (Step~5); formal definitions are in
Section~\ref{sec:formal}.}
\label{fig:architecture}
\end{figure*}
